\scanpage{65}
\[
T_{pC}(M)=T_p(M)\otimes_{\mathbb Q_p}C\simeq\sum T_{\mathrm{Hdg}}(M)^{i,j}\otimes C(-i)
\]
Donc pour tout isom
\[
\deleted{\unsure{Q}}\qquad C(1)\overset{\rho}{\simeq}C
\]
on déduit un isom bigradué
\[
\widetilde\rho:T_{pC}(M)\simeq T_{\mathrm{Hdg}\,C}
\]
i.e. une trivialisation
\[
\widetilde\rho\in P(C)=\operatorname{Hom}_{\operatorname{Spec}L}(\operatorname{Spec}C,P)
\]
D'autre part, si $\rho'$ est un autre isomorphisme, on aura
\[
\rho'=u\rho,\qquad u\in C^*
\]
et on devra avoir
\[
\widetilde{u\rho}=i_1(u^{-1})
\]
Donc, on a \deleted{\unsure{on}} un hom de \deleted{torseurs} \inserted{$C^*$-espaces \uuline{à opérations} \uuline{de} $\pi$}
\diagram{\[\boxed{\begin{tikzcd}
C(1)^*\arrow[r]&\unsure{P(C)}
\end{tikzcd}}\]}
\margin{\unsure{éléments non nuls}}
quand on regarde $P(C)$ comme espace où $C^*$ opère grâce à $i_1^{-1}$; en d'autres termes, on a un hom de $Gm_C$-schémas
\inserted{\diagram{\[
\begin{tikzcd}
\unsure{\check{\mathbf V}|\mathbb V}(C(-1))^*\arrow[r,no head]\arrow[d,equal]&P_C\\
\unsure{\check{\mathbf V}|\mathbb V}(\mathbb Z_p(-1))_C^*
\end{tikzcd}
\]}}
L'hom précédent doit être compatible avec $(j_1,j_2)$ d'une part, $(i_1,i_2)$ d'autre part, i.e. $(j_1,j_2)$ résulte par \uline{descente} de $i_1,i_2$, quand on choisit l'isom.
\[
G_L^P\otimes_L C\simeq G_C
\]
associé à un $\widetilde\rho$ .....
\mspage{64}
